\clearpage
\chapter*{강연 232의 정오표}
\addcontentsline{toc}{chapter}{강연 232의 정오표}
\setcounter{footnote}{9}
\src{303}

\noindent\textbf{7. TDTE V: 피카르 스킴. 존재 정리.}

\begin{center}
[부르바키 세미나, 제14권, 1961/62, 제232호, 19쪽.]
\end{center}

\noindent\textbf{7쪽, 비고 3.3.} --- 여기서 제기한 문제는 멈퍼드에
의해 긍정적으로 해결되었다(다음 강연에 대한 논평을 보라).

\noindent\textbf{13쪽, 비고 5.2.} --- 다음 강연의 첫머리에서
지적하듯이, 여기서 제기한 존재 추측은 거짓이다. 그러나 멈퍼드는 자신의
방법으로 이보다 약간 약한 정리를 증명하는 데 성공했다.

\noindent\textbf{14쪽, 7행.} --- ``완비'' 뒤에 ``잉여체가 대수적으로
닫힌''이라는 조건을 덧붙여야 한다. 곧 ``완비 국소환''을 ``잉여체가
대수적으로 닫힌 완비 국소환''으로 읽어야 한다. 앞서 언급한 멈퍼드의
반례는 실제로 이 제한이 필수적임을 보여 준다.

\noindent\textbf{17쪽, 비고 6.6.} --- 다음 강연의 첫머리에서
지적하듯이, 여기서 제기한 문제에 Murre가 방금 긍정적인 해답을 주었다.

\noindent\textbf{18쪽, 비고 6.7.} --- 여기서 제기한 문제는
긍정적으로 해결된다(다음 강연에 대한 논평의 마지막 문단을 보라).

\noindent\textbf{18쪽, 비고 6.8, 4행.} --- ``정칙'' 대신
``$k$ 위에서 매끄러운''으로 읽어야 한다.

% FGA Canon print-preservation apparatus: atomic integration R1.
\par\smallskip
\begingroup\footnotesize
\noindent\textit{FGA-CANON-ERR-0032. 인쇄본의 독법은 C095에 계속 기록되어 있으며, FGA-CANON-DISP-0034가 현행 본문의 최소 수정을 규율한다.}\par
\endgroup
